% Options for packages loaded elsewhere
\PassOptionsToPackage{unicode}{hyperref}
\PassOptionsToPackage{hyphens}{url}
%
\documentclass[
  11pt,
]{article}
\usepackage{amsmath,amssymb}
\usepackage{iftex}
\ifPDFTeX
  \usepackage[T1]{fontenc}
  \usepackage[utf8]{inputenc}
  \usepackage{textcomp} % provide euro and other symbols
\else % if luatex or xetex
  \usepackage{unicode-math} % this also loads fontspec
  \defaultfontfeatures{Scale=MatchLowercase}
  \defaultfontfeatures[\rmfamily]{Ligatures=TeX,Scale=1}
\fi
\usepackage{lmodern}
\ifPDFTeX\else
  % xetex/luatex font selection
\fi
% Use upquote if available, for straight quotes in verbatim environments
\IfFileExists{upquote.sty}{\usepackage{upquote}}{}
\IfFileExists{microtype.sty}{% use microtype if available
  \usepackage[]{microtype}
  \UseMicrotypeSet[protrusion]{basicmath} % disable protrusion for tt fonts
}{}
\makeatletter
\@ifundefined{KOMAClassName}{% if non-KOMA class
  \IfFileExists{parskip.sty}{%
    \usepackage{parskip}
  }{% else
    \setlength{\parindent}{0pt}
    \setlength{\parskip}{6pt plus 2pt minus 1pt}}
}{% if KOMA class
  \KOMAoptions{parskip=half}}
\makeatother
\usepackage{xcolor}
\usepackage[margin=25mm]{geometry}
\usepackage{color}
\usepackage{fancyvrb}
\newcommand{\VerbBar}{|}
\newcommand{\VERB}{\Verb[commandchars=\\\{\}]}
\DefineVerbatimEnvironment{Highlighting}{Verbatim}{commandchars=\\\{\}}
% Add ',fontsize=\small' for more characters per line
\newenvironment{Shaded}{}{}
\newcommand{\AlertTok}[1]{\textcolor[rgb]{1.00,0.00,0.00}{\textbf{#1}}}
\newcommand{\AnnotationTok}[1]{\textcolor[rgb]{0.38,0.63,0.69}{\textbf{\textit{#1}}}}
\newcommand{\AttributeTok}[1]{\textcolor[rgb]{0.49,0.56,0.16}{#1}}
\newcommand{\BaseNTok}[1]{\textcolor[rgb]{0.25,0.63,0.44}{#1}}
\newcommand{\BuiltInTok}[1]{\textcolor[rgb]{0.00,0.50,0.00}{#1}}
\newcommand{\CharTok}[1]{\textcolor[rgb]{0.25,0.44,0.63}{#1}}
\newcommand{\CommentTok}[1]{\textcolor[rgb]{0.38,0.63,0.69}{\textit{#1}}}
\newcommand{\CommentVarTok}[1]{\textcolor[rgb]{0.38,0.63,0.69}{\textbf{\textit{#1}}}}
\newcommand{\ConstantTok}[1]{\textcolor[rgb]{0.53,0.00,0.00}{#1}}
\newcommand{\ControlFlowTok}[1]{\textcolor[rgb]{0.00,0.44,0.13}{\textbf{#1}}}
\newcommand{\DataTypeTok}[1]{\textcolor[rgb]{0.56,0.13,0.00}{#1}}
\newcommand{\DecValTok}[1]{\textcolor[rgb]{0.25,0.63,0.44}{#1}}
\newcommand{\DocumentationTok}[1]{\textcolor[rgb]{0.73,0.13,0.13}{\textit{#1}}}
\newcommand{\ErrorTok}[1]{\textcolor[rgb]{1.00,0.00,0.00}{\textbf{#1}}}
\newcommand{\ExtensionTok}[1]{#1}
\newcommand{\FloatTok}[1]{\textcolor[rgb]{0.25,0.63,0.44}{#1}}
\newcommand{\FunctionTok}[1]{\textcolor[rgb]{0.02,0.16,0.49}{#1}}
\newcommand{\ImportTok}[1]{\textcolor[rgb]{0.00,0.50,0.00}{\textbf{#1}}}
\newcommand{\InformationTok}[1]{\textcolor[rgb]{0.38,0.63,0.69}{\textbf{\textit{#1}}}}
\newcommand{\KeywordTok}[1]{\textcolor[rgb]{0.00,0.44,0.13}{\textbf{#1}}}
\newcommand{\NormalTok}[1]{#1}
\newcommand{\OperatorTok}[1]{\textcolor[rgb]{0.40,0.40,0.40}{#1}}
\newcommand{\OtherTok}[1]{\textcolor[rgb]{0.00,0.44,0.13}{#1}}
\newcommand{\PreprocessorTok}[1]{\textcolor[rgb]{0.74,0.48,0.00}{#1}}
\newcommand{\RegionMarkerTok}[1]{#1}
\newcommand{\SpecialCharTok}[1]{\textcolor[rgb]{0.25,0.44,0.63}{#1}}
\newcommand{\SpecialStringTok}[1]{\textcolor[rgb]{0.73,0.40,0.53}{#1}}
\newcommand{\StringTok}[1]{\textcolor[rgb]{0.25,0.44,0.63}{#1}}
\newcommand{\VariableTok}[1]{\textcolor[rgb]{0.10,0.09,0.49}{#1}}
\newcommand{\VerbatimStringTok}[1]{\textcolor[rgb]{0.25,0.44,0.63}{#1}}
\newcommand{\WarningTok}[1]{\textcolor[rgb]{0.38,0.63,0.69}{\textbf{\textit{#1}}}}
\setlength{\emergencystretch}{3em} % prevent overfull lines
\providecommand{\tightlist}{%
  \setlength{\itemsep}{0pt}\setlength{\parskip}{0pt}}
\setcounter{secnumdepth}{5}
\ifLuaTeX
  \usepackage{selnolig}  % disable illegal ligatures
\fi
\usepackage{bookmark}
\IfFileExists{xurl.sty}{\usepackage{xurl}}{} % add URL line breaks if available
\urlstyle{same}
\hypersetup{
  pdftitle={The finite-energy equation on all non-Archimedean algebraic curves},
  pdfauthor={P16.4 continuation 269 --- unverified research draft},
  hidelinks,
  pdfcreator={LaTeX via pandoc}}

\title{The finite-energy equation on all non-Archimedean algebraic
curves}
\author{P16.4 continuation 269 --- unverified research draft}
\date{7 October 2026}

\begin{document}
\maketitle

\section{Statement, actual objects, and original
scope}\label{statement-actual-objects-and-original-scope}

This paper proves a proposed \textbf{all-genus curve theorem}, relative
to the eight complete caller grants and the ordinary primary theorems
specified below. It also supplies an unconditional application of formal
section generation to the same original line bundle. All new arguments
are unverified drafts. Original P16.4, which includes every dimension,
remains open in the general nontrivially valued higher-dimensional
sector. The submission has \texttt{candidate=false}.

Throughout the original domain, \(K\) is complete for a real-valued
non-Archimedean absolute value, of characteristic and residue
characteristic zero; the valuation may be trivial, discrete or
nondiscrete. The original \(X/K\) is smooth, geometrically integral and
projective, and \(L\in\operatorname{Pic}(X)\otimes\mathbb Q\) is an
actual ample rational bundle. No numerical replacement of \(L\) is
permitted. Put \(T=X^{\rm an}\) and \(V=(L^n)>0\). A metric is psh
precisely when it is an all-point decreasing net of algebraic weighted
Fubini--Study metrics on this fixed rational bundle, in BE Definition
7.2. We exclude the identically \(-\infty\) metric; other poles remain
literal values. Scalars are relative to a continuous psh reference
\(\theta\).

For continuous scalars the genuine operator and raw energy are

\[
 \operatorname{MA}(v_1,\ldots,v_n),\qquad
 E_\theta(v)=\frac1{n+1}\sum_{r=0}^n
 \int_T v\,\operatorname{MA}(v^{[r]},0^{[n-r]}).
\]

For actual singular endpoints define \(E_\theta(u)\) as the infimum of
continuous psh majorant energies and
\(\mathcal E^1=\{E_\theta>-\infty\}\). The extended energy is available
before asserting finite energy. The raw target condition is

\[
 \mu\geq0\text{ Radon},\quad \mu(T)=V,\quad
 \mathcal E^1\subset L^1(\mu).                                      \tag{1}
\]

In (1) every actual finite-energy scalar is absolutely integrable;
constants are finite real constants. Define the continuous dual

\[
 J_\theta(\mu)=\sup_{v\in C(T)\cap\operatorname{PSH}}
       \left(E_\theta(v)-\int_Tv\,d\mu\right).                       \tag{2}
\]

The complete M grant makes (1) equivalent to finite (2), at the correct
raw mass, with reference covariance. This paper uses that granted
equivalence, and does not offer another target criterion or moment
estimate.

\textbf{Theorem A (curve equation, new draft).} If \(\dim X=1\), for
every target (1) there is an actual \(u\in\mathcal E^1\) on the same
original \(L\) with genuine \(\operatorname{MA}_\theta(u)=\mu\). Every
actual \(\mathcal E^1\) solution equals \(u+c\) at \textbf{every}
analytic point for one finite real \(c\). This includes poles and
type-I/type-IV endpoints after algebraic-closure pullback. The solution
attains the energy-minus-target-integral dual. There is no genus,
invariance, skeleton-support, metrizability, countability or
discreteness restriction. The genuine operator is the one granted by O,
constructed independently of the target. No inverse measure construction
replaces it.

Eight and only eight caller grants are mathematical premises: O217,
D214, N220, T171, Q222, C218, G227 and M257, with all unchanged
clarifications. They are locally registered by note n1. O supplies
genuine mixed operators, all-test integrability, extended-energy
decreasing continuity before E1, weighted E1 decreasing-net laws, and
reference/field laws. D supplies the compact orbit quotient, unique
invariant Radon lift with equal continuous dual, and descent of an
\textbf{already invariant actual} singular psh/E1 metric. C supplies
strict locality and comparison. T settles all original trivially valued
fields/all dimensions; G settles genus zero, including nonsplit conics.
Neither is reproved here. N's norming-atom hypothesis remains
conditional. M supplies finite-U coercivity and supporting attainment
\textbf{only for a known output}; neither M nor D grants inverse
existence. Bodies 250/265 and working256--268 are data to be
reconstructed, not extra theorems.

\section{Actual marked horizontal
bundles}\label{actual-marked-horizontal-bundles}

Work first over a complete algebraically closed nontrivially real-valued
field \(C\) in the retained characteristic domain. Set \(R=C^\circ\),
\(\mathfrak m=C^{\circ\circ}\), \(\kappa=R/\mathfrak m\) and
\(\Gamma=-\log|C^\times|\). Choose \(\varpi\in R\) with
\(0<|\varpi|<1\). Rank one makes the \(\varpi\)-adic and valuation
topologies equivalent; \(R\) is complete and separated. Root extraction
makes the nonzero group \(\Gamma\) divisible and dense. For
\(R_j=R/\varpi^jR\), \(\sqrt{\varpi^jR}=\mathfrak m\), so
\(\operatorname{Spec}R_j\) has one point. Its maximal ideal is nil,
possibly with no uniform nilpotence exponent.

Let \(a=\deg L>0\). Choose an integer \(d\) with the original \(M=dL\)
integral very ample and a nonzero original section \(S_0\). Its divisor
is \(D_0=\sum n_pp\), with finitely many \(C\)-rational points and
\(\sum n_p=da\). The identification \(M=\mathcal O_X(D_0)\) is the one
specified by \(S_0\). If we began over \(K\), take \(d,S_0\) over \(K\)
first. BPR v1 4.22 and 3.13 furnish a marked semistable vertex set
refining any given finite type-II set. Subdivide each self-loop as in
3.14. This retains the full graph, all cycles and lengths. The actual
model is the canonical-affinoid gluing of BPR 4.11/4.15, with generic
fibre the given \(X^{\rm an}\). No formal-carrier algebraization is
guessed. Its components are smooth proper residue curves; the marks
reduce distinctly and smoothly by 4.8.

Here is the complete horizontal Cartier application. At a smooth
residual mark, BPR 4.4 gives a formal étale chart to \(xy=b\). One
coordinate is a unit there, so the chart on that open is étale over the
multiplicative line, hence over an open of the relative affine line with
coordinate \(t\). For an actual affine chart \(\operatorname{Spf}B\)
containing the mark, the actual point \(p\) evaluates every powerbounded
element in \(R\): \(\epsilon:B\to R\). This is a continuous retraction
on constants. Put \(c=t(p)\). The residual étale fibre at \(\bar c\) is
a finite reduced set; \(\kappa\) is algebraically closed. Delete the
other fibre points by a principal localization whose value at \(p\) has
norm one. The residual zero fibre of \(t-c\) is now exactly \(\bar p\).

FK I.5.3.10--11 makes \(B_j=B/\varpi^jB\) étale over the corresponding
localization of \(R_j[t]\) for \textbf{every} \(j\). Thus
\(Q_j=B_j/(t-c)\) is affine étale over \(R_j\) and has the given
evaluation section. Its extended maximal ideal is nil: each element
involves finitely many nil coefficients, which generate a nilpotent
ideal. Hence \(\operatorname{Spec}Q_j\) has exactly the one residual
point. A section of an étale affine morphism is an open and closed
immersion (Stacks 024T, the base change of its diagonal). Its image
contains that point, so it is an isomorphism, including structure
sheaves. Therefore \(Q_j=R_j\) with precisely the original evaluation.

The monic polynomial \(t-c\) is a nonzerodivisor on \(R_j[t]\),
including when \(R_j\) has zero divisors, and remains so after
localization and flat étale base change. If \(h\) evaluates to zero in
\(B\), there is consequently a \textbf{unique} \(q_j\) with
\(h=(t-c)q_j\) modulo \(\varpi^j\). Uniqueness makes these quotients
compatible. Completeness and separation give \(h=(t-c)q\) in \(B\).
Conversely every such multiple evaluates to zero. Thus

\[
 \ker\epsilon=(t-c),\quad B/(t-c)=R,
 \quad t-c\text{ is regular in }B.                                 \tag{3}
\]

The kernel is closed. Its generic quotient is the original point, and
its residual intersection has length one. On overlaps its generators
define the same evaluation ideal, so regularity makes their ratios
units. Off the residual mark the equation is a unit: invert its unit
residue first modulo \(\varpi\), using nilness, then lift by a
convergent geometric series. These local ideals glue to an effective
horizontal Cartier divisor \(H_p\). No unknown ideal is assumed to
commute with completion. The finite sum \(H=\sum n_pH_p\) defines an
actual formal bundle \(\mathfrak M_H\) with generic bundle
\textbf{exactly} \(dL\), and

\[
 A_v=\deg(\mathfrak M_H|_{C_v})
     =\sum_{p\rightsquigarrow C_v}n_p\geq0,
 \qquad\sum_v A_v=da.                                             \tag{4}
\]

\section{Shared heights, actual vertical Cartier factors, and
positivity}\label{shared-heights-actual-vertical-cartier-factors-and-positivity}

On the full connected metric graph \(G\), define \(\Delta_{\rm out}\) by
sums of outgoing slopes, counting both loop incidences. For a vertex
function affine on edges,

\[
 -\sum_vh(v)\Delta_{\rm out}h(v)
      =\sum_{e=vw}\frac{(h(w)-h(v))^2}{\ell_e}.                     \tag{5}
\]

Thus the symmetric matrix has kernel constants and image the zero-sum
space. Choose positive \(Q_v\) of sum \(da\) and solve \(\Delta h=Q-A\).
In the edgeless one-vertex case the bundle is already positive.

Approximate the \textbf{shared global} heights \(h(v)\) by
\(g_v\in\Gamma\). On an oriented edge, set \(s=(g_w-g_v)/\ell\). Choose
rational \(p<s<q\) arbitrarily close to the original slope. Insert

\[
 t=\frac{q\ell-(g_w-g_v)}{q-p}\in\Gamma,\qquad 0<t<\ell,           \tag{6}
\]

and use slope \(p\) before \(t\), \(q\) after. The exact endpoint
heights agree on every incident edge; the bend degree is \(q-p>0\).
Loops use the same endpoint height and strict brackets around zero. With
\(Q_* =\min Q_v\) and \(m=\max(1,\max\operatorname{valence}v)\),
outgoing-slope errors smaller than \(Q_*/(4m)\) preserve every old
degree above \(3Q_*/4\). Every \textbf{new} component has positive bend
degree. Height error plus edge length times slope error controls the
uniform error. No later zero-degree vertex is inserted into a claimed
positive final model.

The points (6) are actual type-II points. Refine the actual model by BPR
4.11. Clear slopes by an integer \(N\) and choose \(c_v\) with
\(-\log|c_v|=Ng_v\). On an actual oriented node chart \(xy=b_e\), with
\(x\) a unit near the \(v\)-component and \(y\) near \(w\), take Cartier
factors

\[
 f_v=c_v,\qquad f_e=c_vx^{m_e},\quad m_e=N\,\text{slope}(g).
\]

On the overlaps the ratios are \(x^{m_e}\) and
\((c_vb_e^{m_e}/c_w)y^{-m_e}\). The scalar ratio has norm one by the
\textbf{exact} endpoint relation. It is an integral unit, not
necessarily equal to one. Flatness and \(xy=b_e\ne0\) make these
legitimate regular meromorphic factors, including negative exponents.
The factors glue since all overlaps compare to the shared component
frame. Cycles impose unit-ratio gluing, not equality of arbitrarily
chosen scalars.

The canonical generic section of \(\mathcal O(V_g)\) in frame \(f^{-1}\)
has norm \(|f|\), so its weight is \(-\log|f|=Ng\circ r_G\) at every
point. At the \(v\)-branch, \(e_v=x^{m_e}e_e\) gives order \(+m_e\), and
the other branch gives \(-m_e\). Hence

\[
 \deg\bigl(\mathcal O(V_g)|_{C_v}\bigr)=N\Delta_{
 \rm out}g(v).                                                     \tag{7}
\]

Now \(\mathfrak A=\mathfrak M_H^N\otimes\mathcal O(V_g)\) has generic
bundle \(dNL\) and positive degree on every component. Their sum is
\(Nda\). Its full metric on \(L\) divides the weight by \(dN\); the
scalar change from the horizontal metric is \(g/d\). More generally, if
a reference bundle represents \(DL\), a scalar perturbation \(h\)
requires vertical weight \(DN h\), and final bundle power \(DN L\). All
subsequent global-section powers multiply this \textbf{entire}
denominator.

\section{Unconditional whole-modulo recognition on the original
bundle}\label{unconditional-whole-modulo-recognition-on-the-original-bundle}

\textbf{Theorem B (new recognition application).} Let \(\mathfrak F\) be
an actual proper admissible topologically finitely presented
\(R\)-formal model with given generic identification with the already
algebraic original \(X^{\rm an}\) and proper nodal field fibre of
dimension one. Let \(\mathfrak A\) be invertible with specified generic
identification \(\mathfrak A_\eta=dNL\). If its degree on every reduced
residual component is positive, its normalized formal metric is an
actual zero-weight algebraic FS metric on the same original \(L\), with
denominator \(dNk\) for some integer \(k>0\). In particular the
preceding semistable construction applies unconditionally. No H1 or
generation antecedent is left on \(R/\varpi\).

\textbf{Proof.} The reductions \(X_j=\mathfrak F\bmod\varpi^j\) are flat
finitely presented schemes over \(R_j\). Restricted series reduce to
polynomials and the defining ideal is finitely generated. Properness is
part of formal properness; alternatively Stacks 0BPG(8) transfers
properness from the actual field fibre through the exact nil-thickening
square, with the locally-finite-type hypotheses verified. No global
nilpotence exponent is used. By Stacks 33.44.15, positive degrees on all
components make \(\mathfrak A_\kappa\) ample. Its immediate
finite-surjective argument uses only proper \textbf{field} schemes.
Apply the general Stacks 37.50.3 (0D2S) to the actual proper
\(X_1\to\operatorname{Spec}R_1\) and actual \(\mathfrak A_1\). The ample
neighbourhood of its sole point is the entire base. Thus
\(\mathfrak A_1\) is ample on the \textbf{whole} \(X_1\). Noetherian
30.21.4 is used only inside the proof of this general theorem, after
finite-stage approximation.

Write \(R_1=\varinjlim_i B_i\) for its finite-type
\(\mathbb Z\)-subalgebras. These rings are Noetherian, their affine
spectra form an inverse system with affine transitions, and all
base-changed schemes have affine transition maps. Stacks 01ZM descends
the actual finitely presented \(X_1\): finite affine presentations,
overlaps, morphisms and eventual equations/cocycles descend. Stacks 0B8W
descends the invertible bundle and its inverse with their evaluation
isomorphism. Stacks 081F and 04AI descend properness and locally-fp
flatness at later indices. Stacks 09MT descends ampleness. In applying
its proof, take an actual finite affine nonvanishing cover by sections
of possibly different positive powers; raise them to a common multiple
\textbf{before} descending them. The loci do not change. The descended
nonvanishing loci eventually are affine and cover, so the common power
is ample, hence so is the original bundle. Choose \textbf{one later
index} where every property holds. This gives a proper flat fp
\(Z/B_i\), an ample invertible \(B\), and exact base changes
\((X_1,\mathfrak A_1)\).

The complete Stacks 01Q3 gives global generation of \(B^k\) for all
sufficiently large \(k\), and a finite generating section list by
quasi-compactness. Stacks 0B5T gives finite \(H^0\) and \(H^q(Z,B^k)=0\)
for every \(q>0\) and sufficiently large \(k\). Only finitely many
degrees occur for one finite affine Čech cover, so choose a
\textbf{single} \(k\) satisfying all of them and generation. These
Noetherian statements are applied to \(Z\), never directly to \(R_1\).

For precision, arbitrary base change here follows from an elementary
universal exactness argument. Choose a finite affine cover of separated
\(Z\); all finite intersections are affine. The nonnegative bounded Čech
complex \(C^0\to\cdots\to C^r\) of \(B^k\) has flat \(B_i\)-modules as
terms: the rings are base-flat and the invertible modules are projective
over those rings. It has no higher cohomology. Starting with
\(Z^r=C^r\), descend through

\[
 0\longrightarrow Z^q\longrightarrow C^q
    \longrightarrow Z^{q+1}\longrightarrow0 \quad(0\leq q<r).
\]

Each flat quotient makes this universally exact and makes its kernel
flat. At degree zero \(Z^0=H^0\). Tensoring with \textbf{any}
\(B_i\)-algebra, including the possibly nonflat \(R_1\), preserves the
entire exactness and \(H^0\). The tensor complex is the actual
base-changed affine Čech complex. Hence

\[
 H^q(X_1,\mathfrak A_1^k)=0\ (q>0),\qquad
 H^0(X_1,\mathfrak A_1^k)=H^0(Z,B^k)\otimes_{B_i}R_1.             \tag{8}
\]

The finite evaluation list generates after pullback by right exactness.
This proof specifically requires a complex concentrated in nonnegative
degrees; a general perfect-complex slogan would not suffice.

Flatness and invertible local frames give exact sheaf sequences on the
common underlying space

\[
 0\longrightarrow\mathfrak A_1^k
 \xrightarrow{\,\varpi^j\,}\mathfrak A_{j+1}^k
 \longrightarrow\mathfrak A_j^k\longrightarrow0.                  \tag{9}
\]

The first arrow identifies the layer by torsion-freeness over \(R\).
Equation (8) makes all transition maps on \(H^0\) surjective. Lift the
\textbf{same finite list} recursively through every \(j\). Invertible
formal sheaves are the inverse limits of these reductions in complete
local frames; the global equalizer of local sections commutes with
inverse limits. We obtain actual formal global sections
\(s_1,\ldots,s_m\) of \(\mathfrak A^k\).

On an affine integral frame \(e\), write \(s_\ell=f_\ell e\). Generation
modulo \(\varpi\) gives \(1=\sum b_\ell f_\ell+\varpi h\) with integral
coefficients. The convergent integral series \((1-\varpi h)^{-1}\) gives
an actual identity \(1=\sum c_\ell f_\ell\). Every integral element has
norm at most one at every point of its generic chart, and therefore

\[
                 \max_\ell|f_\ell(x)|=1\quad\text{for every }x.    \tag{10}
\]

The field comparison is on the \textbf{already original algebraic}
\(X\), not an algebraization of \(\mathfrak F\). FK v5 0.9.2.7 and
I.2.1.4(2) place complete \(R\) in the t.u. adhesive domain. Set
\(U=\operatorname{Spec}R\setminus V(\varpi)=\operatorname{Spec}C\). The
original \(X\to U\) is proper fp and the original \(dNkL\) is coherent.
The printed adhesive case of II.9.4.1--2 identifies its algebraic and
rigid \(H^0\). Berkovich1993 whole §1.6 and Proposition 1.3.4 identify
this with the same good strict Hausdorff projective Berkovich carrier,
using the same affine-frame cocycles. Thus the lifted sections are
actual \(S_\ell\in H^0(X,dNkL)\). They have no common algebraic zero,
since over \(C\) a nonempty common-zero locus has a classical point,
contradicting (10). In an original frame the identity is

\[
          \phi(x)=\frac1{dNk}\max_\ell\log|S_\ell(x)|             \tag{11}
\]

at \textbf{every} point. This is exactly the required original-bundle FS
metric. \(\square\)

The immediate RG premise inside the printed FK proof is checked
separately. An \(\varpi\)-torsion-free finite
\(R\langle T\rangle\)-module is \(R\)-torsion-free by rank one, hence
\(R\)-flat since \(R\) is a valuation ring. Its quotient is flat over
\textbf{\(R/\varpi\)}, not claimed flat over \((R/\varpi)[T]\). The
polynomial morphism is locally fp, the module finite, and the one-point
base has finite weak assassin. RG1971 3.4.6 applies on its entire flat
locus and gives finite presentation. The four original scan pages were
visually read; ordinary journal extraction contains no article text. Its
immediate proof uses fp flat approximation and weak-assassin transfer.
Those recursive foundations and all FK coherent/GFGA foundations have
not been freshly reproved. We use the exact printed theorem in its
checked domain, not the broader merely rigid-Noetherian case announced
for a later reference, and not the omitted formal algebraization of FK
I.10.3.2.

\section{Canonical residue fields and genuine raw formal
curvature}\label{canonical-residue-fields-and-genuine-raw-formal-curvature}

We need the genuine analytic coefficient, not merely a graph degree. On
each actual BPR vertex chart the canonical reduction is an affine open
of the \textbf{same} component \(C_v\). Its Shilov boundary is the
singleton actual type-II point \(v\) (BPR 4.12--4.15). In an edgeless
model take the complement of one open ball as in 4.15.3. A type-II
completed residue field has torsion value-group quotient over \(C\);
divisibility of \(\Gamma\) in its fixed real embedding implies
\(|\mathcal H(v)^\times|=|C^\times|\). Thus every nonzero spectral norm
on this reduced chart lies in the ground value group.

Here is the whole field argument. If a reduced strict affinoid algebra
\(B\) has integral canonical reduction and all nonzero spectral norms in
\(|C^\times|\), normalize nonzero \(f,g\) by scalars of those exact
norms. Their nonzero reductions have nonzero product. Hence the spectral
norm is multiplicative and is an actual point with zero kernel. Every
point reducing generically agrees with it, since normalized nonzero
reductions remain nonzero at the generic point. Finally
\(\mathcal H(v)\) is the completion of the honest
\(\operatorname{Frac}B\). Approximate any powerbounded \(h\) by \(f/g\)
with error \(<1\). Divide both numerator and denominator by a scalar of
\(\|g\|\). The denominator has nonzero canonical residue and the
numerator is powerbounded. Their residue quotient gives the residue of
\(h\). This proves the \textbf{whole} field identity

\[
                   \widetilde{\mathcal H(v)}=\kappa(C_v).          \tag{12}
\]

A singleton reduction fibre alone would not prove (12).

For any actual integral formal invertible bundle on such a model, CLD v1
6.5.1 and 6.6.8 identify its germ residual bundle with its actual
component restriction. The original smooth analytic curve is good,
separated, pure one-dimensional and boundaryless. At nonvertices the
reduction closure is zero-dimensional, so 6.9.2 kills curvature there.
At vertices apply the \textbf{AC branch} 6.9.4 with its entire
calibrated-cone proof 6.9.5 and the value-index remark 6.9.6. Equation
(12) makes the proper component a model of the whole residual germ. The
resulting first-Chern coefficient is its actual degree, without an
unexplained field factor. The formula also holds for signed degrees and
zero degrees: 6.9.2--3 reduce by a difference of globally generated
bundles on the projective residual definition model; the
defining-dimension-zero case has both sides zero. Thus for generic
bundle \(DL\) and normalized metric \(\phi\),

\[
                 c_1(L,\phi)=\frac1D\sum_v
                       \deg(\mathfrak A|_{C_v})\delta_v.           \tag{13}
\]

Local PL approachability follows from CLD 6.8.2--3 and that residual
difference construction. When \(\phi\) is an actual FS metric, its
global smooth semipositive approximations in BE 7.14 identify this CLD
measure with the genuine continuous operator used in O. In dimension one
curvature is linear on signed differences. We do not infer global
approachability of an arbitrary formal metric from local
approachability, or invoke CLD 6.4.4 without its stated hypothesis. For
the positive pairs above (11) supplies global FS recognition first.
Their mass is \(Nda/(dN)=a\) by (7), and curvature is their genuine O
output. The v1 exponent omissions and its general stable-field descent
paragraph are not extrapolated beyond this AC curve application.

Fix one such positive pair and pure FS reference \(\theta_*\). Let
\(\Gamma_0\) be its final finite skeleton and
\(\mu_0=\operatorname{MA}_{\theta_*}(0)\). By (13) it has raw mass
\(a\), is supported on its type-II vertices, and has positive weight on
every component. Fix a vertex \(\xi\) with \(\mu_0(\{\xi\})>0\).

\section{Actual cofinal skeleta and uniform retraction
tests}\label{actual-cofinal-skeleta-and-uniform-retraction-tests}

Order all actual finite semistable skeleta containing \(\Gamma_0\) by
enlargement; subdivide loops if necessary. Write \(\Gamma_j\) and
\(r_j:T\to\Gamma_j\). BPR 3.13 supplies a common refinement and
compatible retractions. Outside \(\Gamma_0\) the complement is a union
of actual open balls, so enlargement only subdivides the core and
attaches finite trees. Every original cycle remains. BPR 5.1 puts every
type-II/III point on some such graph. The whole inverse-limit theorem
5.2 identifies \textbf{all} \(T\), including type-I/IV points, with the
inverse limit of these graphs. No countable cofinal subset is used.

For clarity the resulting uniform convergence is proved without a
metrizability assumption. Continuous functions factoring through one
finite graph form a lattice containing constants and separating points.
For fixed \(x,y\), a separating graph function followed by an affine
rescaling takes prescribed values \(f(x),f(y)\); if these values agree
use a constant. Given \(x\), finite minima of these functions, using a
finite neighbourhood cover in \(y\), give a factor function at most
\(f+\epsilon\) globally, equal to \(f\) at \(x\) and larger than
\(f-\epsilon\) near \(x\). Finite maxima using compactness in \(x\) then
give a factor function \(h\) with \(\|h-f\|_\infty<\epsilon\). All
finitely many factors use one larger graph. Beyond it \(h\circ r_j=h\),
whence

\[
                 \|f\circ r_j-f\|_\infty\longrightarrow0.         \tag{14}
\]

Consequently \(r_{j*}\mu\to\mu\) weakly for every finite Radon measure,
with arbitrary support. This is approximation by a cofinal family, not
restriction of the target to one skeleton.

\section{Grounded kernels retaining every
cycle}\label{grounded-kernels-retaining-every-cycle}

On each compact connected finite metric graph define \(j_\xi(x,y)\) by

\[
 j_\xi(\xi,y)=0,\qquad
 \Delta_{{\rm out},x}j_\xi(x,y)=\delta_\xi-\delta_y.                \tag{15}
\]

Insert \(x,y,\xi\) as vertices. The negative outgoing matrix has
conductances \(1/\ell_e\); its grounded principal minor is positive
definite by (5). Its inverse is symmetric. Discrete minimum principle
gives nonnegative values for its unit source. Linear interpolation gives
the entire solution, including the source kink. Thus (15) has a unique
symmetric nonnegative kernel. It is bounded and jointly continuous:
subdividing an edge at a moving point gives a finite electrical network,
series subdivision preserves the original network, and the linear
solution converges when a subedge shrinks to zero by contracting its
endpoints. Equivalently the identity
\(2j_\xi(x,y)=R(x,\xi)+R(y,\xi)-R(x,y)\) follows by polarization of the
grounded inverse; effective resistance is continuous under moving
subdivision points by the same series law. Loops use both incidences, or
subdivision into parallel edges, and are never deleted.

If a finite tree is attached at \(b\), its new kernel is the old kernel
at the two retractions plus the length of the common initial segment of
the two paths from \(b\) \textbf{when both points are in that same
tree}. The extra term is zero for different attached trees. At \(b\) its
added outgoing unit cancels the old source at \(b\), at \(y\) it gives
the negative unit, and all other outgoing sums cancel; it is grounded.
Uniqueness proves the formula. Iterating handles every enlargement.

Pull back to \(T^2\):

\[
             j_j(x,y)=j_\xi^{\Gamma_j}(r_jx,r_jy).
\]

Each is bounded jointly continuous and the net increases under
enlargement. Its supremum can be infinite, for example on a classical
diagonal. We never integrate that supremum before establishing the
needed energy bound. Since \(\mu_0\) is on \(\Gamma_0\),

\[
 J_0(x)=\int_T j_j(x,v)d\mu_0(v)
\]

is independent of \(j\), bounded continuous and factors through
\(\Gamma_0\). Put \(I_{00}=\int J_0d\mu_0\).

\section{Realizing graph potentials as actual original-bundle
metrics}\label{realizing-graph-potentials-as-actual-original-bundle-metrics}

\textbf{Lemma C (actual finite inverse).} On a fixed actual
\(\Gamma_j\), every positive Radon \(\nu\) of mass \(a\) has a
continuous actual \(\theta_*\)-psh solution

\[
 h_\nu(x)=-\int_{\Gamma_j}j_\xi(r_jx,y)d\nu(y)+J_0(x),
 \qquad \operatorname{MA}_{\theta_*}(h_\nu)=\nu,\quad h_\nu(\xi)=0. \tag{16}
\]

The measure is viewed as an actual measure on \(T\). No assertion that
\(h_\nu\) is PL is made for nonatomic \(\nu\).

\textbf{Proof.} First suppose \(\nu\) is finite atomic on actual type-II
points. Insert them as vertices. Mix a vanishing positive proportion of
a mass-\(a\) measure positive on \textbf{every} vertex, including
zero-reference new vertices. It suffices initially to treat these
strictly positive vertex targets. The unique anchored graph solve is
\(\Delta h=\nu-\mu_0\). The reference is represented by an actual formal
bundle with generic \(DL\). Repeat (6) with shared \(\Gamma\) heights
and rational upward bends applied to \(Dh\), not to an independently
polarized graph. Old component margins are \(D\nu(v)>0\). Clear rational
slopes by \(N\). Twist the pulled-back actual reference bundle to power
\(N\) by the actual vertical Cartier data. It has generic \(DN L\),
positive degree on every old/new component, and its scalar is a uniform
approximation to \(h\) on the \textbf{entire} \(T\). Theorem B on this
particular pair gives genuine FS metrics with full denominator \(DNk\);
(12)--(13) give their genuine raw curvature. Old weights converge to the
desired weights; the finitely many new bend weights tend to zero. Hence
the genuine measures converge weakly to \(\nu\).

Choose a sequence of errors \(e_m\leq2^{-m}\) and add scalar shifts
\(4\,2^{-m}\). The shift drop \(2\,2^{-m}\) exceeds \(e_m+e_{m+1}\).
These genuine weighted FS metrics decrease uniformly to \(h\). This
constant operation is allowed on the original fixed bundle, whether or
not the shift lies in \(\Gamma\). The limit is a literal FS endpoint;
compact boundedness gives finite majorant energy. O's genuine decreasing
law identifies its MA with \(\nu\). Let the mixing proportion decrease
to zero, use joint kernel continuity, and perform the same uniform FS
extraction. This removes strictness of the finite target.

For an arbitrary \(\nu\) on this graph, partition the compact graph into
finitely many small Borel pieces and move each positive mass to a nearby
type-II point. Such points are dense along every edge since \(\Gamma\)
is dense. The measures retain mass \(a\) and their potentials converge
uniformly by joint kernel uniform continuity: the error is at most \(a\)
times the kernel oscillation, uniformly in the first variable. Each
atomic potential just constructed is continuous psh and has uniform FS
approximations. A diagonal choice with summable errors and the same
decreasing shifts gives one genuine FS sequence to \(h_\nu\). Uniform
curvature continuity, or O's bounded decreasing law, now yields (16).
All applications of O occur \textbf{after} actual psh recognition.
\(\square\)

\section{A decreasing genuine inverse for arbitrary intrinsic
targets}\label{a-decreasing-genuine-inverse-for-arbitrary-intrinsic-targets}

Let \(\mu\) satisfy (1) and put \(\mu_j=r_{j*}\mu\), which can be
nonatomic. By Lemma C the scalars

\[
 u_j(x)=-\int_Tj_j(x,y)d\mu(y)+J_0(x)                              \tag{17}
\]

are actual continuous psh and have genuine
\(\operatorname{MA}(u_j)=\mu_j\). They are anchored at zero, bounded
above by \(J_0\), and decrease at \textbf{every} point. Write

\[
 I_j(\nu,\eta)=\iint_{T^2}j_j(x,y)d\nu(x)d\eta(y),\qquad
 C_0=\int_TJ_0d\mu.
\]

All these integrals on one finite graph are finite. Symmetry and the
genuine dimension-one energy identity of O give

\[
 E_{\theta_*}(u_j)=\tfrac12\int_Tu_j(d\mu_j+d\mu_0)
                  =\frac{I_{00}-I_j(\mu,\mu)}2,                  \tag{18}
\]

\[
 \int_Tu_jd\mu=-I_j(\mu,\mu)+C_0,\qquad
 F_\mu(u_j)=\frac{I_j(\mu,\mu)-2C_0+I_{00}}2.                     \tag{19}
\]

In obtaining (18), the cross terms cancel; in (19) \(u_j\) factors
through \(r_j\). The constants preserve the full raw \(a\) throughout. M
gives finite \(J_{\theta_*}(\mu)\) by bounded reference covariance, not
a presumed moment at the new reference. Since (17) are \textbf{actual}
continuous tests, (19) gives

\[
 I_j(\mu,\mu)\leq2J_{\theta_*}(\mu)+2C_0-I_{00}.                 \tag{20}
\]

In particular (18) has a uniform finite lower bound. Moment-only
coercivity is not used.

Let \(u=\inf_j u_j\). Its anchor is zero, so it is not identically
\(-\infty\). Here is a literal single FS net proving it is in the
prescribed actual cone, without an unproved minimum-of-FS or
net-flattening rule. Let

\[
 \mathcal A=\{\phi\text{ genuine FS scalar}:
       \phi\geq u+\delta\text{ on all }T\text{ for some }\delta>0\}.
\]

Uniform FS approximation of one \(u_j\) and a positive shift gives a
member. Given finitely many prior members, let \(\delta\) be the least
of their positive gaps. The directed open sets
\(\{u_j<\min_\ell\phi_\ell-\delta/2\}\) cover compact \(T\). One index
therefore gives this inequality everywhere. Approximate that continuous
\(u_j\) by FS with error \(<\delta/16\) and add \(\delta/8\). This
produces a member below every prior one and still above \(u\) by a
positive gap. Thus \(\mathcal A\) is downward directed; quotient
identical metrics if an actual partial order is desired. Its infimum is
\(u\), because each \(u_j\) has arbitrarily close FS approximations and
arbitrarily small positive shifts.

For every \(\phi\in\mathcal A\) a similar directed-open-cover argument
gives some \(j\) with \(u_j\leq\phi\) everywhere. Hence \(E(\phi)\) is
at least the lower bound from (18)--(20). O's \textbf{extended} energy
continuity on this literal FS net applies before E1; it proves
\(E(u)>-\infty\). Now and only now invoke O's E1 decreasing-net law on
\(u_j\downarrow u\). Equations (14) and (16) yield

\[
                     \operatorname{MA}_{\theta_*}(u)=\mu.         \tag{21}
\]

This is actual inverse existence, not attainment deduced for an already
supplied output. After (21), M's granted supporting inequality proves
actual dual attainment. No compactness of arbitrary families of
normalized potentials or continuity of arbitrary envelopes was needed
for this inverse.

\section{Every FS scalar has a finite actual carrier and trace
law}\label{every-fs-scalar-has-a-finite-actual-carrier-and-trace-law}

We supply the additional metric application needed for \textbf{every}
solution. Let \(v=\phi-\theta_*\) for any genuine weighted algebraic FS
metric \(\phi\) on the original \(L\). Clear its section degree and the
reference degree together. Choose a nonzero algebraic rational section
of the common bundle as frame. The finitely many ratios of its sections
and reference sections have finite divisors. Refine to the actual marked
punctured skeleton using BPR 3.13/4.11. BPR 5.15, with 2.10/2.13 and the
actual ordinary-node BL85 3.2 calculation, makes their log absolute
values factor through it and affine with integral slopes between
vertices and marks.

At a mark, use an algebraic nonvanishing local \(L\)-frame and a
uniformizing parameter. Each section is a finite order times a unit; a
unit norm is constant in a sufficiently small ball. Since the sections
are basepoint free, at least one section has order zero and finite
weight. Terms of positive order tend to \(-\infty\), so the finite
maximum is constant sufficiently far down the ray. The same holds for
the reference. Truncate the finitely many rays beyond all such switches
at actual type-II points. Hence \(v\) factors through a \textbf{finite
actual} skeleton, is continuous and piecewise affine there, and is
constant on its outer tails. Type-III switches caused by real weights
remain as interior breakpoints, not discarded.

First take weights in \(\Gamma\). Every value of \(v\) at a type-II
vertex lies in a rational multiple of \(\Gamma\), hence in \(\Gamma\),
and every slope is rational after the full common section denominators.
Switch locations then lie in \(\Gamma\); make them actual vertices. Pull
back the actual reference formal bundle, and construct a signed vertical
Cartier divisor by the factors already proved in (7), with multiplier
clearing the \textbf{entire} scalar/reference denominators. This actual
bundle has the same original generic bundle and its metric is
\textbf{exactly} \(\theta_*+v\) at every point. There is no positivity
or FS recognition assertion inferred from the abstract graph here:
\(\phi\) is already an FS metric. The signed raw formula (13) and actual
residual fields (12) give

\[
                    \Delta_{
 \rm out}v=\operatorname{MA}_{\theta_*}(v)-\mu_0.                  \tag{22}
\]

Zero-degree new components are allowed in this \textbf{signed curvature
computation}, unlike Theorem B's positive recognition application. Their
coefficient is zero. This distinction removes a spurious strict-positive
recognition premise for arbitrary FS tests.

For arbitrary real weights, approximate the finite list by \(\Gamma\)
weights uniformly. Bounded weight ranges allow a common sufficiently far
type-II truncation of all rays, because the same order-zero sections
dominate every nearby list on the tails. Curvatures converge weakly by
genuine continuous uniform continuity. To avoid any unjustified
convergence of edge derivatives, first prove the trace representation
below for each \(\Gamma\) list, and pass using the bounded jointly
continuous kernel. This establishes the trace law directly for every
real list, including type-III switches.

On any fixed \(\Gamma_j\) containing \(\Gamma_0\), enlarge to contain
its finite carrier. Summing outgoing slopes in each attached finite tree
cancels all internal incidences. Its total raw curvature appears at the
attachment, while the flat outer tails contribute zero. The restriction
to \(\Gamma_j\) therefore has

\[
 \Delta_{\Gamma_j}(v|_{\Gamma_j})=r_{j*}\operatorname{MA}(v)-\mu_0,
\]

or equivalently, by grounded graph uniqueness,

\[
 v(x)-v(\xi)=-\int_T j_j(x,y)d\operatorname{MA}(v)(y)+J_0(x)
                       \quad(x\in\Gamma_j).                       \tag{23}
\]

Finite raw mass and jointly continuous bounded kernels pass (23) under
the uniform \(\Gamma\)-weight approximation, proving it for all genuine
real-weight FS metrics. This proof uses actual Cartier models and
genuine CLD coefficients, not Thuillier curvature named as a substitute
measure.

\section{Every solution agrees at every
endpoint}\label{every-solution-agrees-at-every-endpoint}

Let \(v\) be \textbf{any} actual E1 solution of (21), and take its
defining decreasing weighted-FS net \(v_i\downarrow v\). The anchor is
finite: O's all-test integrability against the genuine output \(\mu_0\),
together with \(\mu_0(\{\xi\})>0\), gives \(v(\xi)>-\infty\).
Boundedness above gives the other bound. This uses E1 of an already
specified solution; it makes no assertion that every E1 potential is
finite everywhere and does not pretend M's old reference atom is an atom
of a new reference. If desired M's before-E1 moment may be used in its
\textbf{own} reference, but is unnecessary here.

O gives weak \(\operatorname{MA}(v_i)\to\mu\). On a fixed finite
\(\Gamma_j\), the kernel family \(\{j_j(x,\cdot):x\in\Gamma_j\}\) is
compact in \(C(T)\) for the uniform norm. A finite \(\epsilon\)-net and
fixed raw mass reduce uniform convergence of its integrals to finitely
many weak tests. Taking the limit in (23) gives

\[
                  v(x)-v(\xi)=u_j(x)\quad(x\in\Gamma_j).           \tag{24}
\]

This proves all-point equality on every finite graph, hence at every
type-II/III point. Endpoints require a further argument; upper
semicontinuity alone is insufficient.

For every continuous FS \(v_i\), on an attached tree outside
\(\Gamma_0\) the reference curvature is zero and its true curvature is
positive. A cut beyond a point collects nonnegative future mass. Summing
(22) on that future subtree with flat tails gives its derivative away
from the core equal to \textbf{minus} that future mass. It is therefore
nonpositive. For real weights this inequality follows also by uniform
passage from the \(\Gamma\) lists. Actual ball retractions follow
ancestor paths. Consequently for every \(x\in T\),

\[
 v_i(r_jx)\geq v_i(x),\qquad
                   \inf_jv_i(r_jx)=v_i(x).                         \tag{25}
\]

The equality has a particularly direct justification: once \(\Gamma_j\)
contains the finite adapted carrier, factorization gives it exactly,
including type-I/IV points in the collapsed outer domains. The
inequality is needed to ensure no smaller earlier value is included in
the infimum.

Now commute the two actual infima, which are infima of the same doubly
indexed set, without a measurable-net convergence assertion. The finite
anchor in (24) gives

\[
 \begin{aligned}
 v(x)&=\inf_i\inf_jv_i(r_jx)
      =\inf_j\inf_iv_i(r_jx)\\
     &=v(\xi)+\inf_ju_j(x)=v(\xi)+u(x)\qquad(x\in T).
 \end{aligned}                                                     \tag{26}
\]

Literal \(-\infty\) values are preserved. This covers poles,
infinite-depth classical endpoints and finite-depth type-IV endpoints.
Neither almost-everywhere equality, Hilbert-space uniqueness nor
comparison alone replaces (25)--(26). Thus \textbf{every} actual
solution is unique modulo a finite constant.

\section{Original-field descent after singular
invariance}\label{original-field-descent-after-singular-invariance}

Let the original nontrivially valued \(K\) be arbitrary in the retained
domain and let \(C=\widehat{\overline K}\). Pull back a continuous
original reference \(\theta\). D supplies the canonical invariant Radon
lift \(\widetilde\mu\), preserving raw mass, with
\(J_{\theta_C}(\widetilde\mu)=J_\theta(\mu)\). M and bounded reference
covariance therefore put the upstairs target in the just-proved domain,
even though the constructed positive \(\theta_*\) and model need not be
invariant. The absolute operator is reference independent by O. Solve
upstairs and normalize its scalar relative to the \textbf{invariant
pulled-back reference} by supremum zero. That supremum is finite because
the actual endpoint is bounded above and nontrivial.

Every automorphic image is a genuine E1 solution for the same invariant
target, by operator naturality and the invariant lift. The
every-solution \textbf{every-point} theorem (26) makes it differ by a
finite constant; supremum normalization makes that constant zero. The
actual singular metric is therefore invariant \textbf{before} descent. D
now descends this all-point psh metric and its energy to the same
original \(L\). O's genuine singular pushforward yields the original
equation with raw mass \(a\), with no field-degree, probability or
factorial factor. No averaging of singular metrics is performed.

For two downstairs solutions, their pullbacks have the same upstairs
target: the operators on invariant slots are invariant by equivariance
and D's unique invariant lift identifies them. Apply (26) upstairs and
surjectivity of the quotient to obtain the constant equality at every
downstairs point. M gives attainment after this genuine inverse.
Trivially valued original fields are already covered by the complete T
grant. This completes the new all-genus/original-field curve
application, without recapping G.

\section{Replacement inputs and remaining original
gaps}\label{replacement-inputs-and-remaining-original-gaps}

The curve proof replaces generic variational compactness by bounded
finite-graph inversion, the actual uniform retraction limit (14),
positive same-bundle FS approximation, the raw finite-dual energy bound
(20), and compact directed-open-cover construction of a literal
decreasing FS net. It uses no unproved global compactness of the whole
normalized curve cone. Envelopes and envelope orthogonality are bypassed
for this existence proof, rather than asserted. Energy differentiation
remains O's genuine affine-segment polynomial law; M's supporting-plane
argument is invoked only after (21). Strict locality and comparison are
C's already granted theorems. The extra every-point uniqueness
replacement is actual adapted-carrier curvature, finite-kernel trace
convergence, radial monotonicity and commuting infima. Original-field
invariance and descent are checked only after every-solution uniqueness.

For \textbf{general \(n\geq2\)} over nontrivially valued nondiscrete
original fields, O gives the genuine operator and C gives
locality/comparison, and M gives the intrinsic target condition and
finite-U coercivity. These do not give an actual maximizing endpoint or
another full inverse. No normalized-potential compactness topology with
adequate closure and upper semicontinuity is established here. The
necessary envelope continuity/orthogonality and both-sign outer energy
differentiation remain missing wherever a variational proof requires
them. The curve's one-dimensional graph Green inverse has no supplied
higher-dimensional substitute. Every-solution every-point uniqueness in
that residual scope is not proved. For nontrivial discrete fields we
assert no new coverage beyond the exact caller grants; no
literature-exhaustiveness claim is made. Neither a full-original closure
nor a new higher-dimensional inverse is claimed.

Material changes from whole265 and working267/268 are:
whole-mod-\(\varpi\) generation and H1 are now derived at one actual
Noetherian stage using the complete new 01Q3 unit; arbitrary base change
is proved by the bounded nonnegative flat Čech complex; original
full-power algebraic FS recognition is applied unconditionally to
positive constructed pairs; signed actual Cartier models avoid
unnecessary positive recognition for already-FS uniqueness tests;
nonatomic graph targets are realized only after actual atomic metric
approximation; the literal decreasing FS net and every-endpoint infimum
proof are supplied; the anchor is justified by O at an already E1
solution, without moving M's reference atom; singular invariance
precedes original-field descent. Whole250 remains failed. No favorable
failed review or newly generated source pointer is promoted to a grant.

\section{Evidence, foundations and rerun
record}\label{evidence-foundations-and-rerun-record}

\texttt{finite\_checks.py} and \texttt{finite-inputs.json} define a
small exact arithmetic domain. They check graph inversion with retained
cycles, shared-height rational bends, positive old/new weights and full
raw denominator, tree-kernel compatibility and trace cancellation, and
raw energy/dual identities at mass \(5/7\). A separate exact interval
example has a nonatomic uniform target and a quadratic, non-PL inverse.
Positive-degree Čech base change is checked on explicit universally
exact complexes, with a negative-degree counterexample recording why the
proof cannot be broadened. These are finite algebra results and
illustrative computations, not verification of the formal geometry or
analytic theorem.

Run the finite checks with:

\begin{Shaded}
\begin{Highlighting}[]
\ExtensionTok{python}\NormalTok{ finite\_checks.py }\OperatorTok{\textgreater{}}\NormalTok{ finite{-}checks.log}
\end{Highlighting}
\end{Shaded}

Build and packet commands are in \texttt{reproduce.sh}. The exact result
is \texttt{finite-output.json}. Input hashes are in
\texttt{input-verification.json}.

The application ledger is \texttt{source-application-ledger.md}. The
program \texttt{verify\_packet.py} extracts supplied bodies and verifies
only the authorized frozen files.

The file \texttt{environment.json} records the existing tools and the
absence of an exposed Xean/native model or supervised compute facility.
No invocation of that service is claimed, no heavy local computation was
launched and no tool/runtime installation was performed. The requested
particular backend cannot be certified from this tool environment.

Primary source statements are ordinary mathematics in their exact
checked domains, not caller grants. We read the immediate BPR
model/refinement, inverse-limit and slope proofs, including their
explicit ordinary-node BL85 application; the complete relevant Stacks
statements/proofs; the literal CLD AC coefficient proof with its entire
calibration lemma; the printed FK adhesive comparison and immediate
affinoid proof; and Berkovich's whole rigid/coherent H0 bridge. The
original RG and BL85 target scan pages were visually checked. Standard
recursive dependencies not freshly reconstructed include stable
reduction, Tate/BGR gluing and acyclicity, EGA/Stacks affine descent and
coherent cohomology foundations, RG flat fp approximation and
weak-assassin lemmas, FK GFGA/I.C.3.3 and Deligne formula, CLD
rational-cone equality and composite-valuation calibration machinery,
compact Hausdorff/Radon/Riesz foundations. Their exact use and retained
hypotheses are disclosed in the ledger. Source hashes, builds, finite
grids and native execution are provenance/evidence, not independent
correctness review or acceptance. A separate
whole-paper/source/application review remains required.

\section{References}\label{references}

\begin{itemize}
\tightlist
\item
  Baker--Payne--Rabinoff, \emph{On the structure of nonarchimedean
  analytic curves}, arXiv:1404.0279v1; especially 2.10--13, 3.13--14,
  4.1--15/22, 5.1--15.
\item
  Bosch--Lütkebohmert, \emph{Stable reduction and uniformization of
  abelian varieties I}, Math. Ann. 270 (1985), 349--379; actual 3.1--2
  node calculation.
\item
  Boucksom--Eriksson, \emph{Spaces of norms, determinant of cohomology
  and Fekete points in non-Archimedean geometry}, arXiv:1805.01016v3;
  fixed-bundle FS/continuous background, with singular theory supplied
  by O.
\item
  Chambert-Loir--Ducros, \emph{Formes différentielles réelles et
  courants sur les espaces de Berkovich}, literal arXiv:1204.6277v1;
  6.5.1, 6.6.8, 6.8.2--3, 6.9.2--6.
\item
  Stacks Project frozen original units 0AYQ/0B5X, 0D2S, 0BPG, 01ZM,
  0B8W, 081F, 04AI, 09MT, 0B5T, 01Q3, 024T, 02GH.
\item
  Fujiwara--Kato, \emph{Foundations of Rigid Geometry I},
  arXiv:1308.4734v5; 0.9.2, I.2.1, I.5.3, II.6.4, II.9.1/3/4, in the
  printed adhesive domain.
\item
  Berkovich, \emph{Étale cohomology for non-Archimedean analytic
  spaces}, Publ. Math. IHÉS 78 (1993); 1.3.4--6 and whole §1.6.
\item
  Raynaud--Gruson, \emph{Critères de platitude et de projectivité},
  Invent. Math. 13 (1971); whole 3.4, printed 23--26.
\item
  Complete supplied caller grants
  O217/D214/N220/T171/Q222/C218/G227/M257, including all unchanged
  clarifications; registered as local note n1.
\end{itemize}

\end{document}
